% Packages, page style and macros, input by main.tex.
\usepackage[letterpaper,hmargin=0.9in,top=0.76in,bottom=0.78in]{geometry}
\usepackage[T1]{fontenc}
\usepackage{lmodern,textcomp}
\usepackage{xcolor}
\usepackage{booktabs,tabularx,array}
\usepackage{xltabular}
\usepackage{flafter}
\usepackage{placeins}
\usepackage{needspace}
\usepackage{capt-of}
\usepackage{newfloat}
\usepackage{listings}
\usepackage{tikz}
\usepackage{tcolorbox}
\tcbuselibrary{breakable}
\usepackage[hidelinks]{hyperref}
\usepackage[hyperfirst=false]{glossaries}

% ============================================================
% References to the main paper: edit here if its numbering changes
% ============================================================
\newcommand{\papertitle}{A Neuro-Symbolic Architecture for Safe and Feasible Self-Recovery in Multi-Agent Manufacturing Systems}
\newcommand{\FigOutline}{Fig.~5}
\newcommand{\FigComposition}{Fig.~6}

% ============================================================
% Page and heading style (as in the template)
% ============================================================
\pagestyle{plain}
\setlength{\parindent}{0pt}
\setlength{\parskip}{4pt}
\raggedbottom
\emergencystretch=2em
\setlength{\textfloatsep}{10pt plus 2pt minus 2pt}
\setlength{\intextsep}{8pt plus 2pt minus 2pt}
\setlength{\abovecaptionskip}{4pt}
\setlength{\belowcaptionskip}{0pt}
\makeatletter
\renewcommand\section{\@startsection{section}{1}{\z@}%
  {14pt plus 4pt minus 2pt}{6pt plus 2pt minus 2pt}{\large\bfseries\boldmath}}
\renewcommand\subsection{\@startsection{subsection}{2}{\z@}%
  {10pt plus 3pt minus 2pt}{3pt plus 1pt minus 1pt}{\normalsize\bfseries}}
\renewcommand\subsubsection{\@startsection{subsubsection}{3}{\z@}%
  {8pt plus 2pt minus 2pt}{2pt plus 1pt minus 1pt}{\normalsize\itshape}}
% On a page that holds only tables, they start at the top with a fixed gap between them.
\setlength{\@fptop}{0pt}
\setlength{\@fpsep}{12pt}
\setlength{\@fpbot}{0pt plus 1fil}
\makeatother

% ============================================================
% Text styles. Two families only: the serif text face, and one monospaced
% face for names and for text copied from the records.
%   body                      12 pt serif
%   tables, captions, key     \footnotesize serif
%   recorded text             9 pt monospaced, in boxes and in table cells
%   labels inside boxes       9 pt serif bold (ours, not recorded)
%   notes on omitted lines    9 pt serif italic, gray (ours)
% Colour is used only for the pass and fail marks, the red bar, and the figure.
% ============================================================
\definecolor{markgreen}{RGB}{56,118,29}
\definecolor{markred}{RGB}{192,0,0}
\definecolor{cutgray}{gray}{0.40}

\lstdefinestyle{record}{
  basicstyle=\fontsize{9}{11}\selectfont\ttfamily,
  columns=fullflexible, keepspaces=true, showstringspaces=false, upquote=true,
  breaklines=true, breakatwhitespace=true, breakindent=9.5pt, breakautoindent=true,
  prebreak={}, postbreak={},
  aboveskip=0pt, belowskip=0pt,
  escapeinside={(*@}{@*)}
}
\lstset{style=record}

% A cut inside a recorded line, and recorded lines left out.
\newcommand{\cut}{{\normalfont\ttfamily\color{cutgray}[...]}}
\newcommand{\vd}{\raisebox{-0.1ex}{\vbox{\baselineskip 3pt\lineskip 0pt\lineskiplimit 0pt\hbox{.}\hbox{.}\hbox{.}}}}
\newcommand{\skipped}[1]{{\normalfont\itshape\fontsize{9}{11}\selectfont\color{cutgray}\vd\enspace #1}}
% A recorded name inside one of our notes or labels: plain monospaced, in the colour of the note.
\newcommand{\secname}[1]{{\normalfont\ttfamily #1}}
% Who wrote the lines that follow (our label, not recorded text).
\newcommand{\who}[1]{{\normalfont\bfseries\fontsize{9}{11}\selectfont\rule{0pt}{2.6ex}#1}}
% A bar in the margin beside the line to notice.
\newcommand{\notebar}{\rlap{\hspace{-5pt}\textcolor{markred}{\rule[-0.35ex]{2pt}{1.9ex}}}}
% Recorded text inside a table cell.
\newcommand{\rec}[1]{{\fontsize{9}{11}\selectfont\ttfamily\frenchspacing #1}}

% Names from the records in running text: the monospaced text face itself,
% so every character comes from one font. A line may break after a slash only.
% They are set slightly smaller than the text around them, so that the two faces look the same size.
\makeatletter
\newcommand{\codesize}{\fontsize{\dimexpr\f@size pt*92/100\relax}{\f@baselineskip}\selectfont}
\DeclareRobustCommand{\code}[1]{{\codesize\ttfamily\frenchspacing\expandafter\code@scan\detokenize{#1}\code@end}}
\def\code@scan#1{\ifx#1\code@end\else\code@put#1\expandafter\code@scan\fi}
\def\code@put#1{#1\ifnum`#1=`\/ \penalty\exhyphenpenalty\fi}
\let\code@end\relax
\makeatother
% A recorded name of several words (an event name), in the same face and size as \code.
\newcommand{\recordname}[1]{{\codesize\ttfamily\frenchspacing #1}}

% Verdict marks, drawn so that they look the same in every engine.
\newcommand{\ok}{\tikz[baseline=-0.55ex]{\draw[line width=0.9pt, color=markgreen!80!black, line cap=round, line join=round] (0,0) -- (0.08,-0.09) -- (0.24,0.12);}}
\newcommand{\no}{\tikz[baseline=-0.55ex]{\draw[line width=0.9pt, color=markred, line cap=round] (0,-0.09) -- (0.19,0.11) (0,0.11) -- (0.19,-0.09);}}
\newcommand{\na}{--}

% ============================================================
% Gray titled boxes (as in the template). A prompt section stands where the text shows it and
% continues on the next page when it does not fit; its numbered caption stays with its end.
% ============================================================
\tcbset{recstyle/.style={colback=black!2, colframe=black!40, colbacktitle=black!10, coltitle=black,
  arc=4pt, boxrule=0.4pt, left=7pt, right=7pt, top=5pt, bottom=5pt, fonttitle=\small\bfseries}}
\newtcolorbox{recbox}[1]{recstyle, breakable, title={#1}, title after break={#1 (continued)},
  before skip=8pt, after={\par\nointerlineskip\nopagebreak}}
% The caption is set in a minipage, so no page break can fall between the box and its caption.
\newcommand{\promptcaption}[2]{\noindent\begin{minipage}{\linewidth}\captionof{promptbox}{#1}\label{#2}%
  \end{minipage}\par\addvspace{10pt}}

\newcolumntype{Y}{>{\raggedright\arraybackslash}X}
% A gray box of recorded text (a prompt excerpt or a reply) is numbered as a Listing, apart from the tables.
\DeclareFloatingEnvironment[name=Listing, placement=htb, within=section]{promptbox}
% Tables that may continue on the next page: caption as wide as the text, little space around.
% (set at the start of the document, once the page geometry is in force)
\AtBeginDocument{\setlength{\LTcapwidth}{\textwidth}}
\makeatletter
% Their captions are laid out like those of the other tables.
\def\LT@makecaption#1#2#3{\LT@mcol\LT@cols c{\hbox to\z@{\hss\parbox[t]\LTcapwidth{\@makecaption{#2}{#3}}\hss}}}
\makeatother
\setlength{\LTpre}{6pt}
\setlength{\LTpost}{2pt}
